\chapter{Public Key Algorithmen}
  Das Konzept der Public-Key-Kryptographie basiert auf der 
  Tatsache, da\3 es m"oglich ist zwei Schl"ussel 
  (einen "offentlichen $E$ zum Verschl"usseln und einen geheimen
  $D$ zum Entschl"usseln) zu benutzen, wobei es aber praktisch 
  unm"oglich ist, den geheimen Schl"ussel aus dem "offentlichen 
  zu berechnen. Viele der vorgeschlagenen Public-Key-Algorithmen 
  stellten sich als unsicher heraus und viele der sicheren Algorithmen sind nicht praktisch 
  einsetzbar, da sie entweder zu gro\3e Schl"ussel ben"otigen oder das Chiffrat
  viel zu gro\3 wird im Vergleich zu der Nachricht. Von den sicheren und praktisch 
  einsetzbaren Algorithmen eignen sich nur einige f"ur eine Schl"usselvergabe. 
  Einige der Algorithmen sind nur zur Verschl"usselung und andere nur zur digitalen 
  Signatur geeignet. Es gibt nur zwei Klassen von Algorithmen, die sich sowohl f"ur 
  Verschl"usselung als auch f"ur digitale Signatur eignen und auch "uber l"angere
  Zeit h"aufig eingesetzt werden. 

  Die eine Klasse basiert auf dem Faktorisierungsproblem; 
  das RSA-Verfahren ist der bekannteste Vertreter. Die andere Klasse basiert 
  auf dem diskreten Logarithmusproblem in Gruppen; hier ist das Verfahren von 
  ElGamal das bekannteste Verfahren. Diese Public-Key-Algorithmen sind 
  langsam im Vergleich zu symmetrischen Verschl"usselungsalgorithmen. 
  Daher werden in heutigen Systemen sowohl Public-Key-Komponenten, als auch 
  symmetrische Komponenten eingesetzt; man spricht von Hybridkryptosystemen.
%% ************** Knapsack
  \section{Knapsack}
     Die Sicherheit der Knapsack-Algorithmen beruht auf der Tatsache, 
     da\3 das allgemeine Knapsackproblem ein NP-vollst"andiges Problem ist. 
     Obwohl sich der urspr"ungliche Algorihtmus sp"ater als unsicher 
     heraustellte, lohnt es sich denoch, ihn zu untersuchen, um zu sehen, 
     wie ein NP-vollst"andiges Problem f"ur die Entwicklung eines 
     Public-Key-Algorithmus verwendet werden kann.
     \begin{definition}[Das Knapsackproblem] Es seien beliebige Gewichte 
       $a_1, \ldots ,a_n\in \N$ und eine beliebige Summe $s \in \N$ gegeben. 
       Es wird eine Teilmenge von $\{ a_1, \ldots ,a_n \}$ gesucht, so da\3 
       sich die in ihr enthaltenen Elemente zu $s$ addieren. Anders formuliert 
       lautet das Problen: Finde ein 
       $(x_1,\ldots ,x_n) \in \{0,1\}^n$ mit $\sum_{i=1}^{n}x_i a_i = s$.
     \end{definition}
     Obwohl das allgemeine Knapsackproblem NP-vollst"andig ist, gibt es 
     spezielle Knaps"acke, die einfach zu l"osen sind. Ein superincreasing 
     Knapsack ist ein Beispiel daf"ur.
     \begin{definition}[Der superincreasing Knapsack]
        Es seien $b_1,\ldots, b_n\in \N$ und $s\in \N$ ein Knapsack, 
        wobei f"ur alle $i$ gilt $\sum_{j=1}^{i-1}b_j < b_i$.
     \end{definition}
     Wegen dieser speziellen Struktur l"a\3t sich der Knapsack 
     $(b_1,\ldots ,b_n)$ mit einem Greedy-Algorithmus von oben her l"osen.

     Das n"achste Beispiel zeigt einen sehr einfachen superincreasing Knapsack.
     \begin{beispiel}[Ein einfaches Knapsackproblem]
        Es seien $a_i = 2^{i-1}$ f"ur $i=1,\ldots , n$. Damit ist es sehr einfach,
        das Knapsackproblem zu l"osen. Der L"osungsvektor $(x_1,\ldots ,x_n)$ 
        entspricht der r"uckw"arts gelesenen Bin"ardarstellung von $s$.
     \end{beispiel}
     Das Knapsackproblem kann wie folgt in einem Public-Key-Algorithmus verwendet werden:
     \begin{itemize}
       \item "offentlicher Schl"ussel: $(a_1,\ldots,a_n) \in \N^n$;
       \item geheimer Schl"ussel: Ein Algorithmus, der den Knapsack 
             $(a_1,\ldots,a_n)$ in einen einfach l"osbaren Knapsack transformiert;
       \item Verschl"usselung: $(x_1, \ldots ,x_n ) \mapsto s=\sum_{i=1}^{n} x_i a_i$;
       \item Entschl"usselung: L"ose das Knapsackproblem $(a_1,\ldots,a_n)$ f"ur die Summe $s$.
     \end{itemize}
     Merkle und Hellman stellten 1976 ein Verfahren vor, das die oben 
     beschriebene Methode benutzt. Sie verwendeten einen superincreasing 
     Knapsack und verschleierten ihn mit einer oder mehreren modularen 
     Multiplikationen.
     \subsection*{Verfahren:}
     \begin{itemize}
        \item W"ahle einen superincreasing Knapsack $(b_1,\ldots ,b_n)$ und 
          verschleiere ihn durch eine lineare Trasformation in einen anderen 
          Knapsack. W"ahle dazu $W,M\in \N$ mit $M > \sum_{i=1}^n b_i$ und 
          $\ggt(M,W)=1$. Berechne 
          \begin{equation} a_i := W b_i \pmod M \label{verschleiern} \end{equation}
          und ver"offentliche das Tupel $(a_1,\ldots,a_n)$ als "offentlichen 
          Schl"ussel. Eine bin"are Nachricht $(x_1,\ldots ,x_n)$ wird verschl"usselt zu 
          $s:=\sum_{i=1}^n x_i a_i$.
        \item Entschl"usselt wird durch Zur"uckf"uhrung auf das urspr"ungliche 
          Problem. Berechne $s' := W^{-1} s \pmod M$ 
          ($W^{-1}$ existiert, da $\ggt(W,M)=1$) und l"ose das 
          Knapsackproblem $(b_1,\ldots,b_n)$ f"ur die Summe $s'$.
     \end{itemize}   
     \subsection*{Angriff auf das Knapsackverfahren:}
     Das Knapsackverfahren weist folgende Schwachstellen auf:
     \begin{itemize}
        \item Es ist linear, d.h. es gilt 
          $\sum_{i=1}^n (x_i + y_i)a_i = \sum_{i=1}^n x_i a_i + \sum_{i=1}^n y_i a_i$ und 
        \item Wenn einige $a_i$ ungerade sind, so ist ein Bit Information 
          "uber den Klartext direkt zu sehen (s gerade oder ungerade).
     \end{itemize}
     Ein echter Angriff basiert auf der Beobachtung, da"s aufgrund der 
     Formel~\ref{verschleiern} Zahlen $k_1, \ldots ,k_n \in \Z$ mit
     \begin{equation} 
       a_i U - k_i M = b_i \label{beobachtung} 
     \end{equation}
     existieren, wobei $U \equiv W^{-1} \pmod M$. Dividiert man beide Seiten der 
     Formel~\ref{beobachtung} durch $a_i M$, so erh"alt man
     \begin{equation} 
       \frac{U}{M} - \frac{k_i}{a_i} = \frac{b_i}{a_i M}. \label{formel} 
     \end{equation}
     Weiterhin gilt $b_i \le 2^{i-n} M$, da die Zweierpotenzen den 
     kleinsten superincreasing Knapsack bilden. Substrahiert man 
     Gleichung~\ref{formel} f"ur $i:=1$ von der
     allgemeinen Gleichung~\ref{formel}, so ergibt sich
     $$ -\frac{k_i}{a_i}+\frac{k_1}{a_1}  =  \frac{b_i}{a_i M} - \frac{b_1}{a_1 M} $$
     Multipliziert man diese Gleichung mit $a_1 a_i$, so erh"alt man
     \begin{equation} 
       -a_1 k_i + a_i k_1 = \frac{a_1 b_i}{M} - \frac{a_i b_1}{M}. 
     \end{equation}
     Es gilt 
     \begin{eqnarray}
       |a_i k_1 - a_1 k_i| & \le & \max \left\{ \frac{a_1 b_i}{M}, \frac{a_i b_1}{M}\right\} \\
                           & \le & 2^{i-n} M.
     \end{eqnarray}
      F"ur sehr kleine $i$ lassen sich damit Ungleichungen in $\Z$ angeben. 
      Dabei ist $a_i$ sehr gro\3, $2^{i-n}$ sehr klein. Es reichen nur wenige 
      Ungleichungen, um die $k_i$ eindeutig festzulegen. Mit dem Algorithmus 
      von Lenstra zum linearen Programmieren lassen sich solche Ungleichungen 
      in polynomialer Zeit l"osen. Im Algorithmus "`Faktorisieren von 
      ganzzahligen Polynomen"' von Lenstra ist ein Teilalgorithmus 
      "`Finden von fast orthogonalen Basen in ganzzahligen Gittern"' enthalten. 
      Dieser Algorithmus wurde 1983 implementiert und konnte auch zum Brechen 
      des Knapsackverfahrens ausgenutzt werden. Eine zweite g"unstige Eigenschaft, 
      die das Verfahren vollst"andig zusammen brechen l"a\3t,ist es, da\3 nicht 
      die urspr"unglichen $U$, $M$ gefunden werden m"ussen. Wenn ein 
      $\frac{u}{m} \approx \frac{U}{M}$ verwendet wird, so entsteht ein 
      "aquivalenter Knapsack, der dieselbe L"osung $(x_1,\ldots, x_n)$ besitzt.
%% ************** RSA (Rivest, Shamir, Adleman)
  \section{RSA-Verfahren}
     Das RSA-Verfahren (benannt nach den Entdeckern Rivest, Shamir, Adleman) ist
     das bekannteste Public-Key-Verfahren. Die Sicherheit von RSA beruht auf der 
     enormen Schwierigkeit, gro"se Zahlen zu faktorisieren. 
     \subsection*{Verfahren:}
     \begin{enumerate}
       \item W"ahle zwei gro"se Primzahlen $p$ und $q$ (70-200 Dezimalstellen oder mehr).
       \item Berechne das Produkt $n = p\cdot q$.
       \item W"ahle einen "offentlichen Schl"ussel $e$ mit $\ggt(e,(p-1)(q-1))=1$.
       \item Berechne den geheimen Schl"ussel $d$ mit 
         $e\cdot d \equiv 1 \pmod{(p-1)(q-1)}$ und 
         ver"offentliche das Tupel $(n,e)$ als 
         "offentlichen Schl"ussel.
     \end{enumerate}
     Mit Hilfe des Chinesische Restsatzes sieht man, da"s 
     $m^{ed} \equiv m \mod{n}$ f"ur alle $m\in\Z_n$ gilt.
     \begin{itemize}
       \item Verschl"usseln: \\
         Es sind der Modulus $n$ und der Schl"ussel $e$ "offentlich bekannt. 
         Zur Verschl"usselung wird eine Nachricht $m$ in $m_1,\ldots ,m_l$ 
         zerlegt mit $m_i < n$, und es werden die 
         Chiffrate $c_i \equiv m_i^e \pmod{n}$ berechnet.
       \item Entschl"usseln: \\
         Berechne $m_i \equiv c_i^d \equiv (m_i^e)^d \equiv m_i^{ed} \equiv m_i^1$ 
         f"ur $i=1,\ldots,l$ und setze daraus die Nachricht $m$ zusammen.
     \end{itemize}
     Das RSA-Verfahren kann auch zum Signieren verwendet werden, 
     indem man mit dem geheimen Schl"ussel $d$ verschl"usselt.
     Die Signatur kann durch Entschl"usseln mit dem "offentlichen
     Schl"ussel "uberpr"uft werden.

     Das RSA-Verfahren ist in Hardware etwa um einen Faktor 1000 und in 
     Software etwa 100 mal langsamer als symmetrische Blockchiffren. 
     Zur Beschleunigung w"ahlt man 
     geeignete einfache Exponenten (es wurden zum Beispiel $e=3, 17$ und 
     $2^{16}+1$ in verschiedenen Standards vorgeschlagen). 
     Es sind durch diese Einschr"ankung bisher keine gr"o"seren Probleme 
     mit der Sicherheit bekannt.
     \subsection*{Primzahlerzeugung:}
     Die Sicherheit des RSA-Verfahrens h"angt in ganz entscheidender Weise 
     davon ab, da\3 die verwendeten Primzahlen geheim bleiben. 
     Deshalb m"ussen beim Erzeugen der Primzahlen nicht vorhersehbare 
     Zahlen generiert werden; sie d"urfen beispielsweise nicht mit einem 
     Zufallsgenerator mit einem anderen Personen bekanntem Anfangszustand
     erzeugt werden.

     Alle Verfahren, gro\3e Primzahlen zu erzeugen, w"ahlen sich zun"achst eine gro\3e 
     Zufallszahl und testen diese auf Primalit"at. Bei den Primzahltests unterscheidet 
     man zwischen zwei Klassen von Algorithmen:
     \begin{itemize}
       \item Probabilistische Primzahltests (relativ schnell und einfach),
       \item Deterministische Primzahltests (recht langsam und aufwendig).
     \end{itemize}
     Die wichtigsten Vertreter der probabilistischen Primzahltests werden
     im n"achten Kapitel vorgestellt.

     \subsubsection*{Angriff auf Protokollebene:}
     Eve belauscht die von Bob an Alice geschickte Nachricht und kennt 
     dadurch das Chiffrat $c \equiv m^e \pmod n$, wobei $e$ der 
     "offentliche und $d$ der geheime Schl"ussel von Alice sind. 
     Eve w"ahlt sich ein $r$ zuf"allig mit $\ggt(r,n)=1$ und 
     berechnet $x \equiv r^e \mod n$ und $y \equiv x c \mod n$. 
     Eve schafft es weiterhin irgendwie, da\3 Alice die Nachricht 
     $y$ signiert (beispielsweise als Zeitstempelsignatur). Damit 
     kennt Eve $y^d$. Sie berechnet 
     $$r^{-1}y^d \equiv r^{-1}x^d c^d \equiv r^{-1}(r^e)^d c^d \equiv r^{-1} r m \equiv m \pmod{n}$$
     und erh"alt damit die Nachricht $m$. Alice f"uhrt eine blinde Signatur aus.

     \subsubsection*{Angriff aufgrund eines gemeinsamen Modulus:}
       Jeder Teilnehmer erh"alt von der Zentrale sein Schl"usselpaar 
       $(e_i, d_i)$, wobei f"ur alle Teilnehmer der Modulus $n$ gleich ist.
       Dadurch entsteht das folgende Problem: eine Nachricht $m$ soll an 
       zwei Personen geschickt werden. Es werden die beiden Chiffrate $m^{e_1}$ 
       und $m^{e_2}$ berechnent und and die entsprechenden Personen 
       geschickt. Es gelte dabei $\ggt(e_1,e_2)=1$ f"ur die geheimen 
       Schl"ussel der Personen. Wenn Eve die beiden Chiffrat $m^{e_1}$ 
       und $m^{e_2}$ belauscht, kann sie 
       $$ r e_1 + s e_2 = 1 $$
       und damit auch 
       $$ (m^{e_1})^r (m^{e_2})^s = m^{r e_1 + s e_2} = m^1$$ 
       berechnen. Deswegen mu\3 f"ur jeden Teilnehmer ein eigenes 
       Paar von Primzahlen $p$ und $q$ gew"ahlt werden.
%% Diffie Hellman
\section{Diffie-Hellman-Verfahren}
  Das Verfahren von Diffie und Hellman erlaubt es, Alice und Bob ein gemeinsames
  Geheimnis "uber eine unsichere Leitung auszutauschen. Das Protokoll kann mit einer
  beliebigen endlichen zyklischen Gruppe realisiert werden. Die Gruppe $G$ und ihr
  Erzeuger $\alpha$ sind "offentlich bekannt.
  \begin{enumerate}
    \item Alice w"ahlt eine Zufallszahl $a$, berechnet $\alpha^a$ in $G$ und sendet
      $\alpha^a$ an Bob.
    \item Bob w"ahlt eine Zufallszahl $b$, berechnet $\alpha^b$ in $G$ und sendet
      $\alpha^b$ an Alice.
    \item Alice empf"angt $\alpha^b$ und berechnet $(\alpha^b)^a$.
    \item Bob empf"angt $\alpha^a$ und berechnet $(\alpha^b)^a$.
  \end{enumerate}
  Alice und Bob haben jetzt ein gemeinsames Gruppenelement $\alpha^{ab}$, das nur ihnen
  bekannt ist. Eve kennt durch Abh"oren nur die Gruppenelemente $\alpha^a$ und 
  $\alpha^b$. Es wird allgemein angenommen, da"s die Gruppenelemente $\alpha^a$ und
  $\alpha^b$ nicht ausreichen, um $\alpha^{ab}$ auszurechnen, wenn es praktisch
  unm"oglich ist, den diskreten Logarithmus von Gruppenelementen zu bestimmen.
  
  Das {\em Diffie-Hellman-Problem} ist es, einen effizienten Algorthmus zu finden,
  der $\alpha^{ab}$ aus $\alpha^a$ und $\alpha^b$ berechnet. Es ist klar, da"s die
  L"osung des diskreten Logarithmusproblem erlaubt, das Diffie-Hellman-Problem zu
  l"osen. Die Umkehrung ist nicht bekannt.

%% ************ ElGamal
  \section{ElGamal-Verfahren}
     Die Sicherheit des Verfahrens von ElGamal
     beruht auf der Schwierigkeit, den diskreten Logarithmus 
     in endlichen zyklischen Gruppen zu bestimmen.
     Es kann sowohl zum Chiffrieren als auch zum Signieren 
     von Nachrichten verwendet werden. 

     Sei $\alpha$ ein Erzeuger der Gruppe $G$. Die Gruppe $G$ und der
     Erzeuger seien "offentlich bekannt. Die Nachrichten werden durch
     Gruppenelemente dargestellt. Alice will Bob eine Nachricht $m$ senden.
     Bob w"ahlt eine Zufallszahl $x$ (sein geheimer Schl"ussel), berechnet
     $y=\alpha^x$ und ver"offentlicht das Gruppenelelement $y$
     (sein "offentlicher Schl"ussel). F"ur den Nachrichtenaustausch
     wird folgendes Protokol verwendet:
     \begin{enumerate}
       \item Alice w"ahlt eine Zufallszahl $k$ und berechnet $a=\alpha^k$.
       \item Alice holt sich Bobs "offentlichen Schl"ussel $y$ und
         berechnet $b=m y^k$.
       \item Alice sendet das Tupel $(a, b)$ an Bob.
     \end{enumerate}
     Bob kann die Nachricht $m$ berechnen, da er seinen geheimen Schl"ussel
     $x$ kennt und $a$ von Alice erh"alt.

     Typischerweise werden die Gruppen $\F_{2^m}^*$, $\F_p^*$ ($p$ Primzahl) oder
     die Gruppe $E(\F)$ der Punkte auf einer elliptischen Kurve "uber einem endlichen
     K"orper benutzt, da sie sich besonders einfach implementieren lassen. Die Gruppe 
     $E(\F)$ ist entweder eine zyklische Gruppe oder
     das direkte Produkt zweier zyklischer Gruppen. Normalerweise ist die eine Komponente des
     direkten Produkts eine kleine zyklische Gruppe und die andere eine gro"se,
     die dann zum Verschl"usseln benutzt wird.
     
     Ein Signaturverfahren l"a"st sich nur mit Primk"orper realisieren. Betrachten
     wir den Fall $G=\F_p^*$ genauer. Bob w"ahlt eine Primzahl $p$, ein primitives 
     Element $g$ und eine Zufallszahl $x$. Er berechnet $y \equiv g^x \pmod{p}$. 
     Sein "offentlicher Schl"ussel ist das Tupel $(y,g,p)$. Sein geheimer Schl"ussel 
     ist der diskrete Logarithmus von $y$ zur Basis $g$, die Zufallszahl $x$.
    
     \subsubsection*{Chiffrierverfahren: }
       Alice verschl"usselt, eine Nachricht $m$, indem sie $k$ zuf"allig mit $\ggt(p-1,k)=1$ 
       w"ahlt und $a \equiv g^k \pmod{p}$ und $b \equiv y^k m$ berechnet. Das Chiffrat
       ist das Tupel (a,b). Bob kann das Chiffrat $(a,b)$ entschl"usselt, indem er
       $$ m \equiv b (y^k)^{-1} \equiv b (g^{xk})^{-1} \equiv b ((g^k)^x)^{-1} \equiv b (a^x)^{-1} \pmod{p}$$
       mit seinem geheimen Schl"ussel $x$ ausrechnet.
     
     \subsubsection*{\bf Signaturverfahren:}
       Bob will eine Nachricht $m$ signieren. Die Idee bei der Signatur 
       ist es, die Nachricht $m$ als
       $$ m \equiv a x \pmod{p-1}$$
       auszudr"ucken. Die Signatur $a$ kann dann durch
       $$ g^{m} \equiv g^{a x} \equiv y^{a} \pmod{p}$$
       "uberpr"uft werden. Der geheime Schl"ussel $x$ kann aber aus der Signatur $a$
       $$ x \equiv m a^{-1} \pmod{p-1} $$
       berechnet werden. Deswegen wird eine lineare Verschiebung benutzt. Dazu 
       w"ahlt Bob $k$ zuf"allig mit $\ggt(k,p-1)=1$ und berechnet 
       $a \equiv g^k \pmod{p}$. Bob berechnet au"serdem noch ein $b$, so da"s
       $$ m \equiv x a + k b \pmod{(p-1)} $$
       gilt. Die Signatur ist dann das Tupel $(a,b)$. Andere Teilnehmer k"onnen 
       Bobs Signatur "uberpr"ufen, indem sie
       $$ g^m \equiv g^{x a + k b} \equiv g^{x a} g^{k b} \equiv y^a a^b \pmod{p} $$
       nachrechnen. Die Signatur funktioniert nur bei Primk"orper,
       weil sonst Probleme mit dem Datentyp auftreten w"urden. Man m"u"st mit einem 
       Gruppenelement potenzieren. Bei Primk"orper werden die Gruppenelmente 
       als nat"urliche Zahlen aufgefa"st. In den anderen F"alle ist dies nicht
       m"oglich.

       Es ist wichtig, da"s jedesmal $k$ verschieden gew"ahlt wird. Falls 
       $a=g^{k} \pmod{p}$ zweimal gleich ist, kann das Linearegleichungssystem
       \begin{eqnarray*}
         m_1 & = & x a + k b_1 \\
         m_2 & = & x a + k b_2 
       \end{eqnarray*}
       aufgestellt und der geheime Schl"ussel $x$ bestimmt werden.
     
%% ******** McElice-Kryptosystem    
   \section{McElice-Kryptosysstem}
     Das McElice-Kryptosystem benutzt fehlerkorrigierende Goppa Codes, die einfach zu 
     decodieren sind, und transformiert sie auf einen "`allgemeinen"' linearen Code. 
     Die Decodierung von allgemeinen linearen Codes ist ein NP-vollst"andiges Problem. 
     Das McElice-System ist folgenderma\3en aufgebaut:
     \begin{itemize}
       \item Geheimer Schl"ussel
         \begin{itemize}
           \item $G'$ ist eine $k\times n$-Matrix, die Generator-Matrix eines Goppa-Codes, 
                 der maximal $t$ Fehler korrigieren kann;
           \item $P$ eine $n\times n$-Permutationsmatrix;
           \item $S$ eine invertierbare $k\times k$-Matrix.
         \end{itemize}
       \item "Offentlicher Schl"ussel
         \begin{itemize}
           \item $G=S G' P$ eine $k\times n$-Matrix;
           \item die Nachricht ist ein Vektor der L"ange $k$ "uber $\{0,1\}$;
           \item verschl"usselt wird durch $c = mG + z$, wobei $z$ ein zuf"alliger Vektor vom Gewicht ein
             wenig kleiner als $t$ ist.
         \end{itemize}
       \item Entschl"usseln
          \begin{itemize}
            \item $c' = c P^{-1}$;
            \item man decodiert $m'$, so da\3 $m'G'=c'$;
            \item es gilt dann $m = m'S^{-1}$.
          \end{itemize}
      \end{itemize}
      Der Vorschlag von McEliece war
      \begin{itemize}
         \item $n = 1024$,
         \item $t \approx 50$ Fehler,
         \item maximal m"ogliches $k$, f"ur das ein Goppa-Code existiert, ist  $k=524$.
      \end{itemize}
      \subsubsection*{Bewertung}
      \begin{itemize}
        \item Der Schl"ussel hat etwa $2^{19}$ bit L"ange.
        \item Die Nachrichtenl"ange wird verdoppelt; der Rechenaufwand ist aber wegen der 
              einfachen Verschl"usselung 
              und Entschl"usselung sehr gering (deutlich schneller als RSA).
        \item Praktisch wird das Verfahren wegen der Schl"usselgr"o\3e kaum eingesetzt.
      \end{itemize}
      \subsubsection*{Angriffe}
      Direktes Decodieren braucht ungef"ahr $2^{500}$ Schritte. Dieser Aufwand l"a\3t 
      sich erheblich reduzieren, wenn man $k$ Spalten der Matrix (mit der Hoffnung, da\3 
      dort kein Fehler dazuaddiert wurde) ausw"ahlt und das lineare Gleichungssystem $mM_k=c_k$ 
      l"ost. Falls kein Fehler in diesen Spalten aufgetreten ist, kann man $m$ 
      durch Matrixinversion zur"uckgewinnen. Es m"ussen allerdings viele Matrizen getestet 
      werden, bis bei 524 Spalten keiner der 50 Fehler enthalten ist. Der Aufwand 
      betr"agt etwa $2^{81}$ Schritte.

 